\documentclass[11pt,a4paper]{article}
\usepackage[margin=2.5cm]{geometry}
\usepackage{amsmath,amssymb}
\usepackage{graphicx}
\usepackage{booktabs}
\usepackage{hyperref}
\usepackage{authblk}
\usepackage{microtype}
\usepackage{enumitem}
\hypersetup{colorlinks=true,linkcolor=blue,citecolor=blue,urlcolor=blue}

\title{Manuscript Title: Concise, Specific, and Informative}
\author[1]{First Author}
\author[2]{Second Author}
\affil[1]{Department, Faculty, University, City, Country}
\affil[2]{Department, Faculty, University, City, Country}
\date{}

\begin{document}
\maketitle

\begin{abstract}
Provide a structured or clearly organized abstract of approximately 200--250 words. State the background/problem, objective, methods, principal results, and conclusion. Do not include citations in the abstract.
\end{abstract}

\noindent\textbf{Keywords:} keyword 1; keyword 2; keyword 3; keyword 4; keyword 5

\section{Introduction}
Define the problem, motivation, research gap, objective, and contributions. Cite prior work using numbered references in order of appearance, for example \cite{sample2026}.

\section{Related Work}
Summarize and critically compare the most relevant literature. Explain how the present study differs from prior work.

\section{Materials and Methods}
Describe datasets, inclusion/exclusion rules, preprocessing, algorithms, software versions, hardware, random seeds, data splits, baselines, hyperparameters, and statistical procedures in enough detail for reproduction.

\subsection{Data and Ethics}
State dataset provenance, licence/permission, ethics approval where applicable, consent requirements, and safeguards for sensitive data.

\section{Results}
Report primary and secondary outcomes with appropriate tables, figures, uncertainty measures, and statistical comparisons where relevant.

\begin{table}[ht]
\centering
\caption{Example results table.}
\begin{tabular}{lcc}
\toprule
Method & Metric 1 & Metric 2 \\
\midrule
Baseline & 0.00 & 0.00 \\
Proposed method & 0.00 & 0.00 \\
\bottomrule
\end{tabular}
\end{table}

\section{Discussion}
Interpret the findings in relation to the research question and prior literature. Separate evidence from speculation.

\section{Limitations}
Describe methodological, data, external-validity, and reproducibility limitations.

\section{Conclusion}
State the main supported conclusion without overstating generalizability.

\section*{Author Contributions}
Use CRediT roles where possible (e.g., Conceptualization, Methodology, Software, Validation, Writing, Supervision).

\section*{Funding}
State funding sources and grant numbers, or: ``This research received no external funding.''

\section*{Conflicts of Interest}
State any competing interests, or: ``The authors declare no conflict of interest.''

\section*{Data and Code Availability}
Provide repository links and access conditions, or explain justified restrictions.

\section*{Ethics Approval and Consent}
Provide the approving body and reference number where applicable, or state why approval was not required.

\section*{Generative AI Disclosure}
Disclose material use of generative AI in writing, coding, image generation, analysis, or translation, including tool, purpose, and human verification. If none, state that no material generative-AI assistance was used.

\begin{thebibliography}{99}
\bibitem{sample2026}
A. Author and B. Author, ``Example article title,'' \textit{Example Journal}, vol. 1, no. 1, pp. 1--10, 2026.
\end{thebibliography}

\end{document}
